% ===== 原文 pp.19–20（PDF pp.20–21）=====
% tr-02 片段 b：(4.5)–(4.15)、4.2 拉氏量主体、4.3 相（开头）
% [原文 p.19]

超对称变换$\delta=i\epsilon_+Q_--i\epsilon_-Q_+-i\overline{\epsilon}_+\overline{Q}_-+i\overline{\epsilon}_-\overline{Q}_+$为
\begin{align}
\delta v_\pm &= \frac{i}{2}\overline{\epsilon}_\pm\lambda_\pm+\frac{i}{2}\epsilon_\pm\overline{\lambda}_\pm,\notag\\
\delta\sigma &= -i\overline{\epsilon}_+\lambda_--i\epsilon_-\overline{\lambda}_+,\notag\\
\delta\lambda_\pm &= i\epsilon_\pm\left((D\pm iv_{01})\pm\frac{1}{2}[\sigma,\overline{\sigma}]\right)+\epsilon_\mp(D_0\pm D_1)\sigma_\mp,\tag{4.5}\\
\delta D &= \frac{1}{2}\epsilon_+\left((D_0-D_1)\overline{\lambda}_++i[\sigma,\overline{\lambda}_-]\right)
+\frac{1}{2}\epsilon_-\left((D_0+D_1)\overline{\lambda}_-+i[\overline{\sigma},\overline{\lambda}_+]\right)+c.c.,\notag
\end{align}
以及
\begin{align}
\delta\phi &= \epsilon_+\psi_--\epsilon_-\psi_+,\notag\\
\delta\psi_\pm &= \pm i\overline{\epsilon}_\mp(D_0\pm D_1)\phi\mp\overline{\epsilon}_\pm\sigma_\mp\phi+\epsilon_\pm F,\tag{4.6}\\
\delta F &= -i\overline{\epsilon}_+(D_0-D_1)\psi_+-i\overline{\epsilon}_-(D_0+D_1)\psi_-
+\overline{\epsilon}_+\overline{\sigma}\psi_-+\overline{\epsilon}_-\sigma\psi_+
+i(\overline{\epsilon}_-\overline{\lambda}_+-\overline{\epsilon}_+\overline{\lambda}_-)\phi.\notag
\end{align}
其中我们仅在本文此处使用记号$\sigma_+=\sigma$与$\sigma_-=\overline{\sigma}$。注意，$\sigma$是A-手征的，而$\phi$是B-手征的：
\begin{equation}
\overline{Q}_+\sigma=Q_-\sigma=0,\qquad \overline{Q}_+\phi=\overline{Q}_-\phi=0.
\tag{4.7}
\end{equation}

% [Tutor Note: p.19 式 (4.7) 与通常教科书的“手征”约定不同，切勿按标准约定去“改正”：本讲义 p.4 给出的定义是A-手征$\overline{Q}_+\cO_A=Q_-\cO_A=0$、B-手征$\overline{Q}_\pm\cO_B=0$（见 PDF p.5 定义那一行），所以 (4.7) 中$\sigma$被$\overline{Q}_+$与$Q_-$消灭、$\phi$被$\overline{Q}_\pm$消灭，两处$\bar Q$的位置恰好与 p.4 的定义逐字一致。核对方法：把 (4.7) 与 p.4 的 A-/B-chiral 定义对照，若照标准约定改写就会与p.4 矛盾。]

在写下拉氏量之前，我们先准备一些记号。我们把$i\fg$上的$G$不变内积与$V$上的$G$不变Hermite内积分别记作
\begin{equation}
(X,Y)\in i\fg\times i\fg\longmapsto\frac{1}{e^2}XY\in\R,\qquad
(\phi_1,\phi_2)\in V\times V\longmapsto\overline{\phi}_1\phi_2\in\C.
\tag{4.8}
\end{equation}

我们还将记$\frac{1}{e^2}XX=\frac{1}{e^2}X^2$、$\overline{\phi}\phi=|\phi|^2$等等。注意，这些内积中参数的个数等于$i\fg$与$V$的不可约分量的个数。例如若$i\fg=\oplus_{k\in K}i\fg_k$是不可约分解，则可写$\frac{1}{e^2}X^2=\sum_{k\in K}\frac{1}{e_k^2}(X_k)^2$，其中$X_k$是$X\in i\fg$的$i\fg_k$分量。常数$e_k$称为第$k$个因子的规范耦合常数。集合$e=(e_k)_{k\in K}$亦可径直称为规范耦合常数。$V$上的内积在$V$上定义一个$G$不变辛结构，我们记$\mu:V\to i\fg^*$为在原点取零的矩映射。

超对称拉氏量为
\begin{align}
\cL &= Q_+Q_-\overline{Q}_+\overline{Q}_-\left(-\frac{1}{2e^2}|\sigma|^2+|\phi|^2\right)
+\Re Q_+Q_-W(\phi)+\Re Q_+\overline{Q}_-(-t(\sigma))\notag\\
&\quad+\text{全导数}\tag{4.9}\\
&=\cL_{\mathrm{g}}+\cL_{\mathrm{m}}+\cL_W+\cL_t,\notag
\end{align}
其中
% [原文 p.20]
% 注：原文中 $\cL_{\mathrm g}$ 的表达式自 p.19 末行跨页延续到本页，编号 (4.10) 印在本页；此处按一条公式排出。
\begin{align}
\cL_{\mathrm{g}} &= \frac{1}{2e^2}\Big[|D_0\sigma|^2-|D_1\sigma|^2+i\overline{\lambda}_-(D_0+D_1)\lambda_-+i\overline{\lambda}_+(D_0-D_1)\lambda_+\notag\\
&\qquad +(v_{01})^2+D^2-\frac{1}{4}[\sigma,\overline{\sigma}]^2-\lambda_+[\sigma,\overline{\lambda}_-]+[\overline{\sigma},\lambda_-]\overline{\lambda}_+\Big],\tag{4.10}
\end{align}
\begin{align}
\cL_{\mathrm{m}} &= |D_0\phi|^2-|D_1\phi|^2+i\overline{\psi}_-(D_0+D_1)\psi_-+i\overline{\psi}_+(D_0-D_1)\psi_+\notag\\
&\qquad +|F|^2+\overline{\phi}D\phi-\frac{1}{2}|\sigma\phi|^2-\frac{1}{2}|\overline{\sigma}\phi|^2-\overline{\psi}_-\sigma\psi_+-\overline{\psi}_+\overline{\sigma}\psi_-\notag\\
&\qquad -i\overline{\phi}\lambda_-\psi_++i\overline{\phi}\lambda_+\psi_-+i\overline{\psi}_+\overline{\lambda}_-\phi-i\overline{\psi}_-\overline{\lambda}_+\phi,\tag{4.11}
\end{align}
\begin{equation}
\cL_W=\Re\left[F^i\frac{\partial W}{\partial\phi^i}-\psi^i_+\psi^j_-\frac{\partial^2W}{\partial\phi^i\partial\phi^j}\right]
\tag{4.12}
\end{equation}
\begin{equation}
\cL_t=-\zeta(D)+\theta(v_{01})
\tag{4.13}
\end{equation}

拉氏量 (4.9) 明显是超对称的，因为$\sigma$是A-手征的而$\phi$是B-手征的（(4.7)）。

场$D$与$F$是“辅助场”——它们没有动能项。它们可以用运动方程消去。这样做之后，我们就得到标量势 (4.14)。

按照约定 (2.13)，theta项进入作用量的方式是
$\frac{1}{2\pi}\int_{\R^2}\dd^2x\,\theta(v_{01})=\int_{\R^2}\theta\big(\frac{i}{2\pi}F_v\big)$，
其中$F_v=iv_{01}\dd x^0\wedge\dd x^1$是取值于$\fg$的（“纯虚”）曲率二形式。这就解释了$\theta$的周期性。

\subsection{相}

标量场的经典势为
\begin{equation}
U(\sigma,\phi)=\frac{1}{8e^2}[\sigma,\overline{\sigma}]^2+\frac{1}{2}|\sigma\phi|^2+\frac{1}{2}|\overline{\sigma}\phi|^2
+\frac{e^2}{2}(\mu(\phi)-\zeta)^2+|\dd W(\phi)|^2.
\tag{4.14}
\end{equation}

注意每一项都非负。$U$的零点集合称为经典真空，它为我们理解理论的低能行为提供了第一条线索。真空方程$U=0$即
\begin{equation}
[\sigma,\overline{\sigma}]=0,\qquad \sigma\phi=\overline{\sigma}\phi=0,\qquad \mu(\phi)=\zeta,\qquad \dd W(\phi)=0.
\tag{4.15}
\end{equation}

后两个方程要求$\phi$位于$\Crit(W)\cap\mu^{-1}(\zeta)$中，前两个方程要求$\sigma$位于$\phi$处稳定子群的嘉当子代数中。FI参数$\zeta$的空间按照$G$-空间$\Crit(W)\cap\mu^{-1}(\zeta)$的拓扑被分割为若干腔室，称为相。在一个相的内部，通常$\Crit(W)\cap\mu^{-1}(\zeta)$的每一点处稳定子群都是有限的，于是$\sigma$被迫取零——经典真空的空间就是商$X_\zeta=(\Crit(W)\cap\mu^{-1}(\zeta))/G$，称为Higgs分支。若是这种情形，理论在低能下约化为Landau–Ginzburg模型$(\mu^{-1}(\zeta)/G,W_\zeta)$，

% [Tutor Note: p.20 式 (4.14) 可由 (4.10)–(4.13) 逐项消去辅助场核出：$D$只出现在 (4.10) 的$\frac{1}{2e^2}D^2$、(4.11) 的$\overline{\phi}D\phi$与 (4.13) 的$-\zeta(D)$中，配方后给出$\frac{e^2}{2}(\mu(\phi)-\zeta)^2$；$F$只出现在 (4.11) 的$|F|^2$与 (4.12) 的$\Re(F^i\partial_iW)$中，$F$的方程给出$F=-\overline{\partial W}$，代回后$|F|^2+\Re(F^i\partial_iW)$在拉氏量中变成$-|\dd W|^2$，故势中出现$|\dd W(\phi)|^2$；而 (4.10) 方括号内的$-\frac{1}{4}[\sigma,\overline{\sigma}]^2$乘以外面的$\frac{1}{2e^2}$，在势中正好是$\frac{1}{8e^2}[\sigma,\overline{\sigma}]^2$。三项系数$\frac{1}{8e^2}$、$\frac{e^2}{2}$可用来反查 (4.10)–(4.13) 的归一化。]
